% ===================================================================
\section{Case Studies}
\label{ch:cases}
% ===================================================================

\subsection{Overview}

Four case studies demonstrate the formalization in operation. Each case study
exercises a different layer of the Lean formalization and illustrates how the
abstract theory connects to concrete execution scenarios.

% -------------------------------------------------------------------
\subsection{Case Study 1: Sepsis Pipeline}
% -------------------------------------------------------------------

\leanref{Example/Sepsis.lean}

\subsubsection*{Problem Statement}

A clinical AI system assists in sepsis diagnosis. The pipeline has three stages:
\begin{enumerate}
  \item \textit{Data aggregation} (Explicit): collects vitals, lab values, EHR
        history.
  \item \textit{Risk scoring} (Explicit): computes qSOFA-based sepsis score.
  \item \textit{Treatment authorization} (Boundary): human physician reviews AI
        recommendation and authorizes treatment.
\end{enumerate}

Tacit-layer components (clinical judgment, patient narrative, prior provider
relationship) are not accessible to the AI at any stage.

\subsubsection*{Formalization}

\begin{lstlisting}
-- Sepsis pipeline definition
def sepsisPipeline : Pipeline where
  stages := [
    ⟨"data_aggregation",   Object.Explicit,  ExplicitOps.aggregate⟩,
    ⟨"risk_scoring",       Object.Explicit,  ExplicitOps.score⟩,
    ⟨"treatment_auth",     Object.Boundary,  BoundaryOps.authorize⟩
  ]
  well_typed := by decide

-- Tacit components are not part of the pipeline
-- (structurally enforced by the type system)
example (t : TacitObject) : ¬ (t ∈ sepsisPipeline.stages.map Stage.obj) := by
  simp [sepsisPipeline]
\end{lstlisting}

\subsubsection*{What the Formalization Verifies}

\begin{enumerate}
  \item The pipeline type-checks: all stages have valid object assignments.
  \item No tacit object appears in the pipeline stage list.
  \item The boundary stage has exactly one well-typed authorization witness.
  \item The pipeline reduces to a unique normal form (by protocol normal form
        theorem).
\end{enumerate}

\subsubsection*{Lean Output}

\begin{lstlisting}[language=bash]
$ lake build Example.Sepsis
[  0 jobs] Compiling Example.Sepsis
[  1 job ] Checking sepsisPipeline : Pipeline
[  1 job ] Checking sepsis_tacit_exclusion : ¬ (...)
[  1 job ] Checking sepsis_auth_uniqueness : ...
Build completed successfully.
\end{lstlisting}

% -------------------------------------------------------------------
\subsection{Case Study 2: Grothendieck Construction}
% -------------------------------------------------------------------

\leanref{Frfp/Core/GrothConstruction.lean}

\subsubsection*{Problem Statement}

The FRFP semantic model requires a category-theoretic construction to integrate
the explicit and tacit categories into a unified semantic universe. The
Grothendieck construction $\int F$ provides this.

\subsubsection*{Key Formal Result}

\begin{lstlisting}
-- The Grothendieck category over the FRFP base category
def frfpGrothendieck : Category := GrothConstruction frfpFunctor

-- Well-formedness: all objects and morphisms satisfy FRFP typing
theorem groth_objects_well_typed :
    ∀ (o : frfpGrothendieck.Obj), WellTypedObj o := by
  intro o
  cases o with
  | mk base fiber => exact well_typed_from_groth base fiber
\end{lstlisting}

\subsubsection*{Proposal P-01 Surfaced Here}

The formalization revealed that the informal paper applies reduction to
categorical \emph{objects}, while the formal construction requires reduction on
\emph{terms} in a term language. This is Proposal P-01
(Chapter~\ref{ch:new}). The case study was necessary to surface this distinction
clearly.

% -------------------------------------------------------------------
\subsection{Case Study 3: The \texorpdfstring{$\mathrm{runTo}\omega$}{runToOmega} Bijection}
% -------------------------------------------------------------------

\leanref{Frfp/Core/RunsTrajectories.lean}

\subsubsection*{Problem Statement}

FRFP traces (sequences of protocol steps) must correspond bijectively to elements
of the semantic domain $\omega$ (the run universe). Establishing this bijection
formally verifies that the operational semantics and denotational semantics are
coherent.

\subsubsection*{Formalization}

\begin{lstlisting}
-- Forward map: trace → ω element
noncomputable def runToOmega : Trace → OmegaElement := fun t =>
  ⟨t.steps.map encode, by exact encode_injective t⟩

-- Backward map: ω element → trace
noncomputable def omegaToRun : OmegaElement → Trace := fun omega =>
  ⟨omega.val.map decode, by exact decode_injective omega⟩

-- Round-trip (bijection proof)
theorem runToOmega_left_inverse :
    ∀ (t : Trace), omegaToRun (runToOmega t) = t := by
  intro t; simp [runToOmega, omegaToRun, encode_decode]

theorem runToOmega_right_inverse :
    ∀ (o : OmegaElement), runToOmega (omegaToRun o) = o := by
  intro o; simp [runToOmega, omegaToRun, decode_encode]
\end{lstlisting}

\subsubsection*{What This Verifies}

The bijection confirms that the FRFP run universe is faithfully modeled. Every
operational trace has a unique denotational representative, and every denotational
element corresponds to a realizable trace. This closes the semantic gap between
the protocol specification and the mathematical model.

% -------------------------------------------------------------------
\subsection{Case Study 4: Semantic Correctness}
% -------------------------------------------------------------------

\leanref{Frfp/Core/SemanticCorrectness.lean}

\subsubsection*{Problem Statement}

The semantic function $\mathrm{Sem} : \mathrm{Term} \to \mathrm{Value}$ must
be proved well-defined (total and type-preserving) and must respect the reduction
relation (reduction-invariant observables).

\subsubsection*{Key Results}

\begin{lstlisting}
-- Sem is well-defined on all well-typed terms
theorem sem_well_defined (t : ProtocolTerm) (ht : WellTyped t) :
    ∃ v : Value, Sem t = v ∧ ValueType v = TermType t := by
  induction t with
  | atom a => exact ⟨semAtom a, rfl, atom_type_preserved a⟩
  | app f arg hf harg =>
      obtain ⟨vf, hvf, htf⟩ := hf
      obtain ⟨va, hva, hta⟩ := harg
      exact ⟨semApp vf va, by simp [Sem, hvf, hva], app_type_preserved htf hta⟩

-- Observable output is invariant under navigation equivalence
theorem observation_invariance (t1 t2 : ProtocolTerm)
    (h_nav : NavEquiv t1 t2) : observe (Sem t1) = observe (Sem t2) := by
  obtain ⟨nf, h1, h2⟩ := nav_equiv_common_nf h_nav
  rw [sem_reduces_to_nf h1, sem_reduces_to_nf h2]
\end{lstlisting}

\subsubsection*{Significance}

Semantic correctness is the bridge between the categorical model and the
executable protocol. These proofs establish that the FRFP implementation (the
explicit computation layer) produces observationally correct outputs regardless
of the internal reduction path taken.
